\section{Diffie-Hellman, ElGamal, dan Pengantar ECC}

\subsection{Diffie-Hellman Key Exchange}

Diffie-Hellman memungkinkan dua pihak untuk \textbf{menyetujui kunci rahasia bersama} melalui kanal tidak aman \cite{stallings2017}. Parameter publik: bilangan prima $p$ dan generator $g$ dari $\mathbb{Z}_p^*$.

\begin{enumerate}
  \item Alice memilih $a$ acak, menghitung $A = g^a \bmod p$, mengirim $A$ ke Bob.
  \item Bob memilih $b$ acak, menghitung $B = g^b \bmod p$, mengirim $B$ ke Alice.
  \item Alice menghitung $K = B^a \bmod p = g^{ab} \bmod p$.
  \item Bob menghitung $K = A^b \bmod p = g^{ab} \bmod p$.
\end{enumerate}

Keduanya memperoleh $K = g^{ab} \bmod p$. Penyerang yang melihat $A$, $B$, $g$, $p$ harus menyelesaikan \textbf{masalah Diffie-Hellman}: hitung $g^{ab}$ dari $g^a$ dan $g^b$, yang setara dengan logaritma diskrit \cite{menezes1996}.

\begin{figure}[htbp]
  \centering
  \begin{tikzpicture}[node distance=1.2cm]
    \node[draw] (alice) {Alice};
    \node[draw, right=3cm of alice] (bob) {Bob};
    \node[below=0.6cm of alice, font=\small] {$a$, $A=g^a$};
    \node[below=0.6cm of bob, font=\small] {$b$, $B=g^b$};
    \draw[->] (alice) -- node[above] {$A$} (bob);
    \draw[->] (bob) -- node[below] {$B$} (alice);
    \node[below=1.8cm of alice, font=\small] {$K=A^b=g^{ab}$};
    \node[below=1.8cm of bob, font=\small] {$K=B^a=g^{ab}$};
  \end{tikzpicture}
  \caption{Diffie-Hellman: pertukaran $A$, $B$; kunci bersama $K=g^{ab}$.}
\end{figure}

\begin{catatan}
Diffie-Hellman sendiri tidak memberikan otentikasi; rentan man-in-the-middle. Digunakan dalam protokol seperti TLS dengan otentikasi sertifikat.
\end{catatan}

\subsection{ElGamal}

ElGamal adalah skema enkripsi berbasis logaritma diskrit \cite{stallings2017}. Parameter: prima $p$, generator $g$.

\textbf{Pembangkitan kunci}: Pilih $x$ acak (privat), hitung $y = g^x \bmod p$ (publik). Kunci publik $(p, g, y)$; kunci privat $x$.

\textbf{Enkripsi} pesan $m$: Pilih $k$ acak, hitung $c_1 = g^k \bmod p$ dan $c_2 = m \cdot y^k \bmod p$. Ciphertext $(c_1, c_2)$.

\textbf{Dekripsi}: $m = c_2 \cdot (c_1^x)^{-1} \bmod p = c_2 \cdot (g^{xk})^{-1} \bmod p = m \cdot y^k \cdot (g^{xk})^{-1} = m$.

ElGamal bersifat probabilistik (bergantung pada $k$ acak). Keamanan bergantung pada kesulitan logaritma diskrit \cite{menezes1996}.

\subsection{Pengantar Elliptic Curve Cryptography (ECC)}

ECC menggunakan grup titik pada kurva eliptik atas field hingga \cite{stallings2017}. Masalah yang mendasari: \textbf{Elliptic Curve Discrete Logarithm Problem (ECDLP)}—diberikan titik $P$ dan $Q = kP$, cari $k$.

Keunggulan: keamanan setara dengan ukuran kunci lebih pendek. Kunci ECC 256-bit setara dengan RSA 3072-bit \cite{nist800_56b_r2}. Digunakan dalam TLS, Bitcoin, dll.

Algoritma ECC: ECDH (pertukaran kunci), ECDSA (tanda tangan), ECIES (enkripsi). Implementasi memerlukan aritmetika kurva eliptik; gunakan library (OpenSSL, libsodium) \cite{menezes1996}.

